% Part: first-order-logic
% Chapter: tableaux
% Section: provability-consistency

\documentclass[../../../include/open-logic-section]{subfiles}

\begin{document}

\iftag{FOL}
      {\olfileid{fol}{tab}{prv}}
      {\olfileid{pl}{tab}{prv}}

\olsection{\usetoken{S}{derivability} و سازگاری}

اکنون شماری از ویژگی‌های رابطهٔ !!{derivability} را اثبات می‌کنیم. این
ویژگی‌ها مستقلاً جالب‌اند، اما هر یک در اثبات قضیهٔ تمامیت نقشی خواهند
داشت.

\begin{prop}\ollabel{prop:provability-contr}
  اگر $\Gamma \Proves !A$ و $\Gamma \cup \{!A\}$ ناسازگار باشد، آنگاه
  $\Gamma$ ناسازگار است.
\end{prop}

\begin{proof}
مجموعه‌های متناهی $\Gamma_0 = \{!B_1, \dots, !B_n\}$ و $\Gamma_1
  =\{!C_1, \dots, !C_m\} \subseteq \Gamma$ وجود دارند به‌طوری‌که
  \begin{align*}
    \{\sFmla{\False}{!A}, &
    \sFmla{\True}{!B_1}, \dots, \sFmla{\True}{!B_n}\} \\
    \{\sFmla{\True}{!A}, &
    \sFmla{\True}{!C_1}, \dots, \sFmla{\True}{!C_m}\}
  \end{align*}
  !!{tableau}s بسته دارند. با استفاده از قاعدهٔ \Cut{} بر $!A$ می‌توانیم
  این دو را در !!{tableau} بسته‌ای یگانه ترکیب کنیم که نشان می‌دهد
  $\Gamma_0 \cup \Gamma_1$ ناسازگار است. چون $\Gamma_0
  \subseteq \Gamma$ و $\Gamma_1 \subseteq \Gamma$، داریم $\Gamma_0 \cup
  \Gamma_1 \subseteq \Gamma$؛ پس $\Gamma$ ناسازگار است.
\end{proof}

\begin{prop}
\ollabel{prop:prov-incons}
$\Gamma \Proves !A$ هرگاه و تنها هرگاه $\Gamma \cup \{\lnot !A\}$
ناسازگار باشد.
\end{prop}

\begin{proof}
نخست فرض کنید $\Gamma \Proves !A$؛ یعنی !!{tableau} بسته‌ای برای
\[
\{\sFmla{\False}{!A},
\sFmla{\True}{!B_1}, \dots, \sFmla{\True}{!B_n}\}
\]
وجود دارد. با استفاده از قاعدهٔ $\TRule{\True}{\lnot}$ می‌توان آن را به
!!{tableau} بسته‌ای برای مجموعهٔ زیر تبدیل کرد:
\[
\{\sFmla{\True}{\lnot !A},
\sFmla{\True}{!B_1}, \dots, \sFmla{\True}{!B_n}\}.
\]

از سوی دیگر، اگر برای مجموعهٔ دوم !!{tableau} بسته‌ای وجود داشته باشد،
می‌توانیم با حذف هر فرمولی که از اعمال \TRule{\True}{\lnot} بر فرض
نخست~$\sFmla{\True}{\lnot !A}$ حاصل شده است، همراه با خود آن فرض، و
افزودن فرض $\sFmla{\False}{!A}$، آن را به !!{tableau} بسته‌ای برای
مجموعهٔ نخست تبدیل کنیم. زیرا اگر شاخه‌ای پیش‌تر از آن رو بسته بود که
نتیجهٔ اعمال \TRule{\True}{\lnot} بر
$\sFmla{\True}{\lnot !A}$، یعنی $\sFmla{\False}{!A}$، را دربر داشت، شاخهٔ
متناظر در !!{tableau} تازه نیز بسته است. اگر شاخه‌ای در تابلوی قدیمی از
آن رو بسته بود که هم فرض $\sFmla{\True}{\lnot !A}$ و هم
$\sFmla{\False}{\lnot !A}$ را دربر داشت، می‌توانیم با اعمال
$\TRule{\False}{\lnot}$ بر $\sFmla{\False}{\lnot !A}$ و به‌دست‌آوردن
$\sFmla{\True}{!A}$، آن را به شاخه‌ای بسته تبدیل کنیم. این کار شاخه را
می‌بندد، زیرا $\sFmla{\False}{!A}$ را به‌عنوان فرض افزوده‌ایم.
\end{proof}

\begin{prob}
ثابت کنید $\Gamma \Proves \lnot !A$ هرگاه و تنها هرگاه
$\Gamma \cup \{!A\}$ ناسازگار باشد.
\end{prob}

\begin{prop}\ollabel{prop:explicit-inc}
  اگر $\Gamma \Proves !A$ و $\lnot !A \in \Gamma$، آنگاه $\Gamma$
  ناسازگار است.
\end{prop}

\begin{proof}
  فرض کنید $\Gamma \Proves !A$ و $\lnot !A \in \Gamma$. آنگاه
  $!B_1$، \dots، $!B_n \in \Gamma$ وجود دارند به‌طوری‌که
  \[
  \{\sFmla{\False}{!A}, \sFmla{\True}{!B_1}, \dots, \sFmla{\True}{!B_n}\}
  \]
  تابلویی بسته دارد. فرض \sFmla{\False}{!A} را با
  \sFmla{\True}{\lnot !A} جایگزین کنید و نتیجهٔ اعمال
  \TRule{\True}{\lnot} بر \sFmla{\True}{\lnot !A} را پس از فرض‌ها درج
  کنید. هر !!{sentence} در !!{tableau} که در !!{tableau} قدیمی با ارجاع به
  سطر~$1$ توجیه می‌شد، اکنون با ارجاع به سطر~$n+2$ توجیه می‌شود. پس اگر
  !!{tableau} قدیمی بسته بود، تابلوی تازه نیز بسته است. این تابلو نشان
  می‌دهد $\Gamma$ ناسازگار است، زیرا همهٔ فرض‌ها در~$\Gamma$ هستند.
\end{proof}

\begin{prop}\ollabel{prop:provability-exhaustive}
  اگر $\Gamma \cup \{!A\}$ و $\Gamma \cup \{\lnot !A\}$ هر دو ناسازگار
  باشند، آنگاه $\Gamma$ ناسازگار است.
\end{prop}

\begin{proof}
  اگر $!B_1$، \dots، $!B_n \in \Gamma$ و $!C_1$، \dots،
  $!C_m \in \Gamma$ وجود داشته باشند به‌طوری‌که
  \begin{align*}
    \{\sFmla{\True}{!A}, &
    \sFmla{\True}{!B_1}, \dots, \sFmla{\True}{!B_n}\} \text{ و}\\
    \{\sFmla{\True}{\lnot !A}, &
    \sFmla{\True}{!C_1}, \dots, \sFmla{\True}{!C_m}\}
  \end{align*}
  هر دو !!{tableau}s بسته داشته باشند، می‌توانیم !!{tableau} ترکیبی و
  یگانه‌ای بسازیم که ناسازگاری $\Gamma$ را نشان دهد. برای این کار
  $\sFmla{\True}{!B_1}$، \dots، $\sFmla{\True}{!B_n}$ را همراه با
  $\sFmla{\True}{!C_1}$، \dots، $\sFmla{\True}{!C_m}$ فرض می‌گیریم و
  سپس قاعدهٔ \Cut{} را اعمال می‌کنیم. حاصل دو شاخه است که یکی با
  $\sFmla{\True}{!A}$ و دیگری با $\sFmla{\False}{!A}$ آغاز می‌شود.

  در سوی چپ، بخش زیر فرض‌های !!{tableau} نخست را بیفزایید. در اینجا هر
  کاربرد قاعده همچنان درست است، زیرا همهٔ فرض‌های !!{tableau} نخست، از
  جمله $\sFmla{\True}{!A}$، در دسترس‌اند. پس هر شاخهٔ زیر
  $\sFmla{\True}{!A}$ بسته می‌شود.

  در سوی راست، بخش زیر فرضِ !!{tableau} دوم را بیفزایید و حاصل هر کاربرد
  ~$\TRule{\True}{\lnot}$ بر $\sFmla{\True}{\lnot !A}$ را حذف کنید.
  نتیجهٔ $\TRule{\True}{\lnot}$ بر $\sFmla{\True}{\lnot !A}$ همان
  $\sFmla{\False}{!A}$ است که با این همه در دسترس است، زیرا نتیجهٔ
  قاعدهٔ \Cut{} در سوی راستِ !!{tableau} ترکیبی است.

  اگر شاخه‌ای در تابلوی دوم از آن رو بسته بود که هم فرض
  $\sFmla{\True}{\lnot !A}$ (که دیگر به‌عنوان فرض در !!{tableau} ترکیبی
  ظاهر نمی‌شود) و هم $\sFmla{\False}{\lnot !A}$ را دربر داشت، می‌توانیم
  $\TRule{\False}{\lnot}$ را بر $\sFmla{\False}{\lnot !A}$ اعمال کنیم تا
  $\sFmla{\True}{!A}$ به دست آید. اکنون شاخهٔ متناظر در !!{tableau}
  ترکیبی نیز بسته می‌شود، زیرا نتیجهٔ سمت راستِ قاعدهٔ \Cut{}، یعنی
  $\sFmla{\False}{!A}$، را دربر دارد. اگر شاخه‌ای در !!{tableau} دوم به
  هر دلیل دیگری بسته شده باشد، شاخهٔ متناظر در !!{tableau} ترکیبی نیز
  بسته می‌شود، زیرا هر !!{signed formula}s جز
  $\sFmla{\True}{\lnot !A}$ که در شاخهٔ !!{tableau} دومِ قدیمی رخ می‌دهد، در
  شاخهٔ متناظرِ !!{tableau} ترکیبی نیز رخ می‌دهد.
  \end{proof}

\end{document}
